\begin{claim}{B.41}{Proved}{The six genus-three actions with Klein-bottle quotient}
There are exactly six mixed finite-subgroup classes in
$\mathcal T_{3,0}$ whose quotient orbifold has underlying
Klein bottle.  Their positive subgroups, full groups, cone orders,
and section counts are
\[
\begin{array}{c|c|c|rrrr}
G^+&G&\text{cone orders}&f&p&|\Phi|&|\Phi/G|\\\hline
C_2&C_4&2,2&0&4&4&2\\
C_2&C_2^2&2,2&0&4&4&2\\
C_3&C_6&3&0&2&2&1\\
C_4&C_8&2&0&4&4&1\\
C_4&C_4\times C_2&2&0&4&4&1\\
C_2^2&D_8&2&0&4&4&1
\end{array}                                             \tag{B.41.1}
\]
Here $f$ and $p$ denote full and nontrivial partial section classes,
and $D_8$ has order eight.  The files
\path{type_3_0_klein_quotient_catalogue.csv} and
\path{type_3_0_klein_quotient_sections.json} contain
every group operation, monodromy, local multiplier, section
stabilizer and centralizer.  All centralizers in the surface
group are trivial.  In the primary positive catalogue, these
rows lie over O3.3, O3.5, O3.7 and O3.8, respectively in
numbers $2,1,2,1$.
\end{claim}
\begin{proof}
The orientation double of a Klein bottle is a torus.
An interior cone of order $m$ in the quotient has two lifts of
order $m$ in that double, and there are no mirrors or corners.
The complete positive catalogue in B.33 therefore leaves
precisely the positive actions
\[
 (C_2;2,2,2,2),\quad(C_3;3,3),\quad
 (C_4;2,2),\quad(C_2^2;2,2),                          \tag{B.41.2}
\]
with torus underlying quotient.
The other torus-positive rows have a single cone and cannot
arise by this double-cover comparison.
Thus the full groups have orders four, six or eight.

We enumerate all oriented group extensions, rather than
presupposing a list of abstract groups.
For a fixed positive group $H$ choose a negative representative
$t$, put $\alpha(h)=tht^{-1}$ and $A=t^2\in H$.
Associativity is equivalent to
\[
           \alpha^2=\operatorname{Inn}(A),
           \qquad\alpha(A)=A.                         \tag{B.41.3}
\]
Conversely the set $H\times\{0,1\}$ with normal form $ht^i$
and relations $th=\alpha(h)t$, $t^2=A$ is a group for every
pair in (B.41.3); checking the two possible negative factors
verifies associativity.  Changing $t$ to $kt$ changes the pair to
\[
 \bigl(\operatorname{Inn}(k)\alpha,\ k\alpha(k)A\bigr),
                                                        \tag{B.41.4}
\]
and changing the identification of $H$ by $\beta\in\operatorname{Aut}(H)$
changes it to
$(\beta\alpha\beta^{-1},\beta(A))$.
These are all orientation-preserving isomorphisms of the
extensions, because any such isomorphism restricts to an
automorphism of $H$ and sends a negative representative to
one negative coset representative.

For $H=C_2$ this gives $C_4$ and $C_2^2$.
For $H=C_3$, the two possible conjugation actions give $C_6$
and $S_3$.  Explicitly the identity action allows the square
of $t$ to be removed using (B.41.4), while inversion forces
$A=1$.
For $H=C_4=\langle a\rangle$, the identity action has two
square classes, $A=1$ or $a$, giving $C_4\times C_2$ and
$C_8$.  Inversion allows $A=1$ or $a^2$, giving $D_8$ and
$Q_8$.
For $H=C_2^2$, the identity action gives $C_2^3$ or
$C_4\times C_2$ according as $A=1$ or $A\ne1$.
Every nonidentity order-two automorphism of $H$ is conjugate
to a transvection.  Its fixed subgroup equals the image of
$1+\alpha$, so (B.41.4) removes either allowable square.
This gives one further oriented extension, $D_8$, with positive
subgroup a four-group.  These arguments prove all the extension
counts $2,2,4,3$ used in the computation.

First take a Klein bottle with one cone of order $m$.
Its orbifold group is
\[
 \Lambda_m=
 \langle d_1,d_2\mid(d_1^2d_2^2)^m=1\rangle.
                                                        \tag{B.41.5}
\]
Put $y=d_1$, $z=d_1d_2$.  Then $y$ is negative, $z$ is
positive, and a positive local cone generator is
\[
                     x=(yzy^{-1}z)^{-1}.              \tag{B.41.6}
\]
A smooth orientation-compatible epimorphism to an oriented
finite group $G$ is therefore exactly a pair $(Y,Z)$ such that
$Y$ is negative, $Z$ is positive, $\langle Y,Z\rangle=G$,
and the image of (B.41.6) has order $m$.
Gromadzki--Szepietowski \cite[Lemma 3.3]{C_GromadzkiSzepietowski2016}
prove that the full outer automorphism group of $\Lambda_m$
is generated by
\[
 (y,z)\longmapsto(yz,z),\qquad
 (y,z)\longmapsto(y^{-1},z),\qquad
 (y,z)\longmapsto(y,z^{-1}).                           \tag{B.41.7}
\]
Their proof identifies these maps with the Dehn twist and
puncture and crosscap slides of the one-pointed Klein bottle.
The automorphism assertion concerns $\Lambda_m$ itself and
does not impose their later nonorientable-kernel hypothesis.
We apply that assertion to the orientation-compatible
epimorphisms above; every generator in (B.41.7) preserves the
orientation character.  Quotienting the finite set of pairs
also by $\operatorname{Aut}(G)$ preserving orientation is
therefore exactly topological conjugacy of the surface actions.

For $H=C_3$ the cone order is three.
In the $S_3$ extension a negative element inverts every
positive element, so (B.41.6) is the identity.
In $C_6$ it equals $Z^{-2}$; $Z$ must generate the positive
$C_3$.  The first move in (B.41.7) runs through the entire
negative coset, and the last move exchanges the two choices
of $Z$.  Thus there is one class, with full group $C_6$.

For $H=C_4$ the cone order is two.
The negative elements of $D_8$ and $Q_8$ invert $H$, again
making (B.41.6) trivial.  In the two abelian extensions it is
$Z^{-2}$, so $Z$ must generate $C_4$.
For $C_4\times C_2$, multiplication of $Y$ by powers of $Z$
sets $Y$ to the order-two generator of the other factor.
For $C_8$, an automorphism sets the odd generator $Y$ to a
chosen generator, and inversion of $Z$ identifies its two
remaining possibilities.  Each gives exactly one class.

For $H=C_2^2$, the elementary abelian extension $C_2^3$
cannot be generated by two elements.
In the abelian extension $C_4\times C_2$, every $Z\in H$
has square one, so (B.41.6) is trivial.
Write the remaining extension as
\[
 D_8=\langle r,s\mid r^4=s^2=1,\ srs=r^{-1}\rangle,
 \qquad H=\langle r^2,s\rangle .
\]
The admissible $Z$ are $s,r^2s$, and every negative $Y$
gives the cone image $r^2$.
An orientation-preserving automorphism exchanges the two
choices of $Z$; the first move in (B.41.7) changes a negative
reflection to a negative rotation, and inversion then
identifies $r$ and $r^{-1}$.  All eight pairs are one orbit.
This gives the final row in (B.41.1).

It remains to handle the two order-two cones in the first case
of (B.41.2).  For either extension $G=C_4$ or $C_2^2$, let
$h$ be the nonidentity positive element.
Both cone generators must map to $h$.
The two crosscap generators may independently map to either
negative element.  Every one of these four assignments is
surjective and respects the relation
\[
                      d_1^2x_1d_2^2x_2=1.
                                                        \tag{B.41.8}
\]
They are all equivalent.  Here is a geometric source move with
its exact word comparison, so no unproved transitivity is needed.
In a one-crosscap neighborhood of a cone, slide that cone once
around the one-sided core.  On the free generators $(d,x)$,
with exterior boundary word $d^2x$, the move is
\[
        d\longmapsto dx,\qquad
        x\longmapsto (dx)^{-1}x^{-1}(dx).              \tag{B.41.9}
\]
It preserves $d^2x$ exactly and sends the cone generator to a
conjugate of its inverse.  It is an automorphism: first replace
the free basis $(d,x)$ by $(dx,x)$, then invert and conjugate
the second member by the first.  Cutting the crosscap
neighborhood along its core gives the stated cone slide;
the unchanged exterior boundary permits extension by the
identity to the other crosscap.
In either abelian target, where $x$ has order two, this
replaces the image of $d$ by its product with $h$ and leaves
the cone image unchanged.  Apply it independently at the two
cones to connect all four assignments.  This gives one class
for each of $C_4$ and $C_2^2$.  Their different exponents
distinguish the two full groups.

For each surviving monodromy the local cone groups inject.
The orbifold covering subgroup is therefore torsion-free,
and it lies in the orientation kernel by construction.
The Klein-bottle orbifold Euler characteristic is
$-\sum_j(1-1/m_j)$; multiplying it by $|G|$ gives $-4$ in
all six cases.  The covering surface is consequently closed,
orientable, and of genus three.  This proves existence and
faithfulness of every listed action, as well as completeness.

There are no mirrors, so no reflection sections.
For each cone generator $x_j$ of prime order $m_j=2$ or $3$,
the points of its critical orbit are exactly
$G/\langle x_j\rangle$.  At the coset represented by $g$,
the stabilizer is $g\langle x_j\rangle g^{-1}$.  Write
$\varepsilon:G\to\{\pm1\}$ for orientation.  The multiplier
map $\mu_{j,g}$ on the section represented by that point is
\[
      \mu_{j,g}(gx_j^kg^{-1})=
                  \varepsilon(g)\frac{k}{m_j}\pmod{\Z}.
                                                        \tag{B.41.10}
\]
Every such point supplies one nontrivial partial section,
which is maximal; its centralizer in the surface group is
trivial by B.3.  No whole $G$ is a local stabilizer, so there
are no full sections.  The resulting coset counts are
exactly (B.41.1).  The program records the permutations of
these cosets, their stabilizers and the multiplier of every
local element.  In the $C_6$ row the two point generators
have multipliers $1/3$ and $-1/3$ under the fixed positive
generator; the negative coset exchanges the two.

The executed standard-library program
\path{catalogue_genus3_klein.py} independently enumerates
all pairs (B.41.3), checks associativity of their finite
multiplication tables and the orientation homomorphism,
enumerates every smooth generating monodromy, and computes
its orbit under the complete source and target automorphism
actions.  It retains the unsuccessful oriented extensions
with their empty monodromy lists, so the exclusions above
are part of the reproducible computation.
\end{proof}
